\documentclass[10pt,a4paper]{amsart}
\usepackage[margin=19mm]{geometry}
\usepackage{amsmath,amssymb,mathtools}
\usepackage[scr=rsfs,cal=euler]{mathalfa}
\usepackage{xfrac}
\usepackage[unicode,hidelinks,bookmarks=false]{hyperref}
\newcommand{\ie}{i.e.\ }
\newcommand{\cf}{cf.\ }
\newcommand{\resp}{respectively\ }
\newcommand{\Subsection}{\subsection}
\providecommand{\nobreakdash}{\nobreak\hbox{-}}
\setlength{\emergencystretch}{2em}
\multlinegap0pt
\usepackage{bm}
\usepackage{bbm}
\usepackage{dsfont}
\usepackage{xspace}
\usepackage{cancel}
\usepackage{mathtools}
\usepackage{amsthm}
\makeatletter
\@ifundefined{Corollary}{\newtheorem{Corollary}{Corollary}}{}
\@ifundefined{Proposition}{\newtheorem{Proposition}{Proposition}}{}
\@ifundefined{Remark}{\newtheorem{Remark}{Remark}}{}
\makeatother
\def\even{{\text{\rm even}}}
\def\odd{{\text{\rm odd}}}
\def\pgfl{\emph{p.g.fl.~}}
\def\spin{{[-1,1]}}
\title[Gaussian Fields on a hypercube from Long Range Random Walks]{Gaussian Fields on a hypercube from Long Range Random Walks}
\author{Robert Griffiths}
\date{Source: arXiv:2510.18167v1 (CC BY 4.0)}
\begin{document}
\maketitle

\noindent\emph{Source and licence.} Gaussian Fields on a hypercube from Long Range Random Walks. Version arXiv:2510.18167v1 (CC BY 4.0), distributed under the Creative Commons Attribution 4.0 International licence, as recorded in the arXiv licence metadata for this exact version and shown on the abstract page https://arxiv.org/abs/2510.18167v1. Full rights notice: licenses/arxiv-2510.18167-CC-BY-4.0.txt.\par\smallskip

\noindent\textit{Editorial scope.} Counted unit: the Corollary whose source label is \texttt{complex:88} in the article (source0.tex lines 1625--1651, source label \texttt{complex:88}), delivered together with the article's own complete original proof of it (source0.tex lines 1653--1717, one \texttt{proof} environment of 65 lines). No mathematics is written, repaired or re-derived: statement and proof are the article's own text line for line. Required context retained uncounted: Gstar:00 (lines 1448--1451); decomposition:22 (lines 1472--1478); prop:complex (lines 1551--1587); Parseval:00 (lines 1580--1585). Every definition, display and label of those lines is retained unchanged. Omitted from this package and not redistributed: 1--1447, 1452--1471, 1479--1550, 1586--1624, 1652--1652, 1718--1789. Nothing inside a retained excerpt is omitted or altered: every retained body is delivered exactly as the article's lines are written, so no omission receipt of any kind accompanies this package. Citations: the retained excerpts carry no citation command at all, so no bibliography entry of the article is transcribed and no reference is redistributed. Reference adaptation: none was needed and none was made; every reference inside the delivered unit resolves inside it, so no unresolved reference or citation is produced. Printed slips kept literal: no printed word, symbol, hypothesis or number of the counted statement or of its proof was altered. Layout and class: the article's own class is \texttt{elsarticle} with packages including natbib, changes, url, bm, amsthm, amsmath, amsfonts, amssymb, tikz, graphicx, enumitem, verbatim; the wrapper is the lane's pinned \texttt{amsart} editorial wrapper at 10pt on A4, which restates the theorem-like environments the retained text needs and adds the article's own macro definitions that the retained text needs (\texttt{\textbackslash even}, \texttt{\textbackslash odd}, \texttt{\textbackslash pgfl}, \texttt{\textbackslash spin}), with extra wrapper packages (bm, bbm, dsfont, xspace, cancel, mathtools). The delivered numbering is the unit's own. Assets: no retained span contains a figure environment, a graphics-inclusion command, a PDF-page inclusion, a table or a TikZ picture; no image file is copied into this package, the package declares no asset dependency and no image file is redistributed with it. Licence: version arXiv:2510.18167v1 of this article is distributed under the Creative Commons Attribution 4.0 International licence, as recorded in the arXiv licence metadata for this exact version and shown on the abstract page https://arxiv.org/abs/2510.18167v1. A full rights notice is delivered at licenses/arxiv-2510.18167-CC-BY-4.0.txt. Characters: every retained excerpt is pure ASCII, so the delivered engine, XeLaTeX with its default Latin Modern text font, reports no missing character.
\begin{equation}
G\big [f\big ]=\frac {1}{1 + \int_\spin (1-f(\xi))\nu_*^\spin(d\xi)}
\label{Gstar:00}
\end{equation}

\begin{align}
\varkappa_t
&= \frac{1}{\sqrt{2\pi}}e^{-\frac{1}{2}t^2}
\sum_{k=0}^\infty \mathbb{E}\big [Y_{\frac{1}{2}}^k\big ]
\frac{1}{\sqrt{k!}}H_k(t)\cdot \zeta_k,
\label{decomposition:22}
\end{align}

\begin{Proposition}\label{prop:complex}
$(\widehat{\varkappa}_\theta)$ has a covariance matrix
\begin{equation*}
\text{\rm Cov}\big (\widehat{\varkappa}_\theta ,\widehat{\varkappa}_\phi\big )
= \mathbb{E}\big [\widehat{\varkappa}_\theta \overline{\widehat{\varkappa}}_\phi\big ]
= e^{-\frac{1}{2}(\theta^2+\phi^2)}\mathbb{E}\big [e^{\theta\phi Y}\big ], 
\label{complexcv:00}
\end{equation*}
where $Y$ is the product of points in the \pgfl (\ref{Gstar:00}).
%\[
%\frac{1}{1+\int_{[-1,1]}\big (1 - f(\xi)\big )\nu_*^\spin(d\xi))}.
%\]
Then
$\widehat{\varkappa}_\theta$ is distributed as $U_\theta+iV_\theta$ where $U_\theta$ and $V_\theta$ are independent real Gaussian processes with 
$U_\theta = \int_{\mathbb{R}_+} e^{i\theta t}\varkappa_t^\even dt$,
$V_\theta = i\int_{\mathbb{R}_+} e^{i\theta t}\varkappa_t^\odd dt$
and
\begin{align}
 \text{\rm Cov}\big (U_\theta,U_\phi\big )
 &= 
 e^{-\frac{1}{2}\big (\theta^2+\phi^2\big )}
 \mathbb{E}\big [\cosh\big (\theta\phi Y\big )\big ]\nonumber \\
  \text{\rm Cov}\big (V_\theta,V_\phi\big )
 &= 
 e^{-\frac{1}{2}\big (\theta^2+\phi^2\big )}
\mathbb{E}\big [ \sinh\big (\theta\phi Y\big )\big ].
\label{covoddeven:00}
\end{align}
Plancherel's theorem shows that
\begin{align}
\mathbb{E}\int_{\mathbb{R}} \lvert\varkappa_s\rvert^2ds &= 
\mathbb{E}\int_{\mathbb{R}} \lvert\widehat{\varkappa}_\theta\rvert^2d\theta\nonumber \\
& = \mathbb{E}\Bigl [\sqrt { \frac{\pi}{1-Y} }\Bigr ],
\label{Parseval:00}
\end{align}
assuming that the expectation on the right of (\ref{Parseval:00}) is finite.
\end{Proposition}

\begin{Corollary}
A strong representation of $U_\theta, V_\theta$ in Proposition \ref{prop:complex} is
\begin{align}
U_\theta &= \sum_{k=0}^\infty
\frac{ \text{\rm Poisson}\big (k;\frac{1}{2}\theta^2\big ) 
 }
{ \sqrt{(2k-1+\delta_{k0})!!}  }
\mathbb{E}\big [Y_{\frac{1}{2}}^{2k}\big ](-\sqrt{2})^k  \zeta_{2k}
\nonumber\\
%%%
V_\theta &= \theta\sum_{k=0}^\infty
\frac{ \text{\rm Poisson}\big (k;\frac{1}{2}\theta^2\big )   }
{ \sqrt{(2k+1)!! } }
\cdot \mathbb{E}\big [Y_{\frac{1}{2}}^{2k+1}\big ](-\sqrt{2})^k \zeta_{2k+1},
\label{complex:88}
\end{align}
where 
$
\text{\rm Poisson}\big (k;\lambda\big ) = \frac{\lambda^k}{k!}e^{-\lambda},\ k \in \mathbb{N}.
$ \\[0.2cm]
An inversion formula is
\begin{equation}
\varkappa_t = \frac{1}{2\pi}\int_{\mathbb{R}}
e^{-i\theta t}\widehat{\varkappa}_\theta\ d\theta,\ t \in \mathbb{R}.
\label{inversion:23}
\end{equation}
\end{Corollary}

\begin{proof}
The representations in (\ref{complex:88}) follow from taking the transform of (\ref{decomposition:22})  and identifying real and imaginary parts. When $X$ is N$(0,1)$,
$
\mathbb{E}\big [e^{i\theta X}H_k(X)\big ]
 = e^{-\frac{1}{2}\theta^2}(i\theta)^k, k \in \mathbb{N}.
$
Taking the transform
\begin{align}
\widehat{\varkappa}_\theta & =\int_{\mathbb{R}}e^{i\theta t}\varkappa_tdt\nonumber\\
& = e^{-\frac{1}{2}\theta^2}
\sum_{k=0}^\infty 
\mathbb{E}\big [Y_{\frac{1}{2}}^k\big ]\frac{1}{\sqrt{k!}}(i\theta)^k\zeta_k\label{Poissonpf:01}\\
&= e^{-\frac{1}{2}\theta^2}\sum_{k=0}^\infty 
\mathbb{E}\big [Y_{\frac{1}{2}}^{2k}\big ]\frac{1}{\sqrt{(2k)!}}(-1)^k\theta^{2k}\zeta_{2k}\nonumber\\
&+ie^{-\frac{1}{2}\theta^2}\sum_{k=0}^\infty 
\mathbb{E}\big [Y_{\frac{1}{2}}^{2k+1}\big ]\frac{1}{\sqrt{(2k+1)!}}(-1)^k\theta^{2k+1}\zeta_{2k+1}.
%\label{Poissonpf:00}
\nonumber
\end{align}
The Poisson expressions now follow using the identities
\begin{align*}
\frac{1}{\sqrt{(2k)!} }= \frac{ \sqrt{2}^k }{2^k}\cdot 
\frac{1}{ \sqrt{(2k-1+\delta_{k0})!!}  }\text{~and~}%\quad
\frac{1}{\sqrt{(2k+1)!}} = \frac{\sqrt{2}^k}{2^k}\cdot \frac{1}{\sqrt{(2k+1)!!}}\ .
\end{align*}
%
To find the inversion formula (\ref{inversion:23}) invert the series (\ref{Poissonpf:01}).
Note that 
\[
\frac{1}{\sqrt{2\pi}}\int_{\mathbb{R}} e^{-i\theta t} e^{-\frac{1}{2}\theta^2}(i\theta)^k d\theta
\]
is the coefficient of $\frac{z^k}{k!}$ in 
\[
\frac{1}{\sqrt{2\pi}}\int_{\mathbb{R}} e^{-i\theta t}
 e^{-\frac{1}{2}\theta^2}e^{i\theta z}d\theta
 = e^{-\frac{1}{2}t^2}\cdot e^{-\frac{1}{2}z^2+zt},
\]
which is $e^{-\frac{1}{2}t^2}H_k(t).$ 
Therefore from (\ref{Poissonpf:01}) and (\ref{decomposition:22})
\[
\frac{1}{2\pi}\int_{\mathbb{R}} e^{-i\theta t}
\widehat{\varkappa}_\theta\ d\theta
= \frac{1}{\sqrt{2\pi}}e^{-\frac{1}{2}t^2}\sum_{k=0}^\infty 
\mathbb{E}\big [Y_{\frac{1}{2}}^k\big ]\frac{1}{\sqrt{k!}}H_k(t)\cdot\zeta_k = \varkappa_t.
\]
\begin{Remark}
Let $W_\theta^u$ and $W_\theta^v$ be independent Poisson$(\frac{1}{2}\theta^2)$ random variables. Then marginally with respect to $U_\theta, V_\theta$
\begin{align}
U_\theta &= 
\mathbb{E}_{W_\theta^u}\Big [
\Big ((2W_\theta^u-1+ I\big \{W_\theta^u=0\big \})!!\Big )^{-\frac{1}{2}}
\mathbb{E}\big [Y_{\frac{1}{2}}^{2W_\theta^u}\big ](-\sqrt{2})^{W_\theta^u } \zeta_{2W_\theta^u}\Big ]
\nonumber\\
%%%
V_\theta &= \theta
\mathbb{E}_{W_\theta^v}\Big [
\Big ((2W_\theta^v+1)!! \Big )^{-\frac{1}{2}}
\cdot \mathbb{E}\big [Y_{\frac{1}{2}}^{2W_\theta^v+1}\big ](-\sqrt{2})^{W_\theta^v} \zeta_{2W_\theta^v+1}\Big ].
%\label{complex:88aa}
\nonumber
\end{align}
These random sums can be extended to a collection $\theta_1,\theta_2,\ldots, \theta_r$ with independent $\big (W_{\theta_j}^u, W_{\theta_j}^v\big )_{j\in [r]}$.
\end{Remark}

\end{proof}

\end{document}
